\documentclass{article}
\usepackage[margin=1in]{geometry}
\usepackage{problemset}
\showsolntrue

% Custom commands
\newcommand{\R}{\mathbb{R}}
\newcommand{\X}{\mathcal{X}}
\newcommand{\Y}{\mathcal{Y}}

\newcommand{\iid}{\stackrel{\text{iid}}{\sim}}
\newcommand{\E}{\mathbb{E}}
\newcommand{\V}{\text{Var}}

\begin{document}

\pstitle{%
    DSC 140A%
}{
    Homework 1%
}[due Monday, October 6 at 11:59PM]

\begin{probset}
\begin{prob}[\textbf{MLE Basics}]
Suppose we are given $n$ data points drawn independently from a Normal distribution:

\[ x_1, \dots, x_n \iid \text{Normal}(\mu, \sigma^2) \]

\begin{subprobset}
    \begin{subprob}
        What is the likelihood function $L(\mu, \sigma^2 | x_1, \dots, x_n)$?

        \begin{soln}
            \begin{align*}
                L(\mu, \sigma^2 | x_1, \dots, x_n) &= \prod_{i=1}^n f(x_i | \mu, \sigma^2) \\
                &= \prod_{i=1}^n \frac{1}{\sqrt{2\pi\sigma^2}} \exp\left(-\frac{(x_i - \mu)^2}{2\sigma^2}\right) \\
                &= \left(\frac{1}{\sqrt{2\pi\sigma^2}}\right)^n \exp\left(-\frac{1}{2\sigma^2} \sum_{i=1}^n (x_i - \mu)^2\right)
            \end{align*}
        \end{soln}
    \end{subprob}

    \begin{subprob}
        What is the maximum likelihood estimator $\hat{\mu}_{MLE}(x_1, \dots, x_n)$?

        \begin{soln}
            \begin{align*}
                \hat{\mu}_{MLE} &= \arg\max_{\mu} L(\mu, \sigma^2 | x_1, \dots, x_n) \\
                &= \arg\max_{\mu} \left[ \left(\frac{1}{\sqrt{2\pi\sigma^2}}\right)^n \exp\left(-\frac{1}{2\sigma^2} \sum_{i=1}^n (x_i - \mu)^2\right) \right] \\
                &= \arg\max_{\mu} \left[ -\frac{1}{2\sigma^2} \sum_{i=1}^n (x_i - \mu)^2 \right]
            \end{align*}
            Since this is a negative sum of squares, the maximum is achieved when 
            \begin{align*}
                \frac{\partial}{\partial \mu} \frac{1}{2\sigma^2} \sum_{i=1}^n (x_i - \hat{\mu}_{MLE})^2 = 0 \\
                \sum_{i=1}^n x_i - n\hat{\mu}_{MLE} = 0 \\
                \hat{\mu}_{MLE} = \frac{1}{n} \sum_{i=1}^n x_i \\
                \implies \hat{\mu}_{MLE} &= \bar{x}
            \end{align*}
        \end{soln}
    \end{subprob}

    \newpage

    \begin{subprob}
        What is the expected value of the estimator $\E (\hat{\mu}_{MLE})$? Is it biased?

        \begin{soln}
            \[
                \E (\hat{\mu}_{MLE}) = \E \left(\frac{1}{n} \sum_{i=1}^n x_i\right) = \frac{1}{n} \sum_{i=1}^n \E (x_i) = \frac{1}{n} \sum_{i=1}^n \mu = \mu
            \]
            Since $\E (\hat{\mu}_{MLE}) = \mu$, $\hat{\mu}_{MLE}$ is unbiased.
        \end{soln}
    \end{subprob}

    \begin{subprob}
        What is the variance of estimator $\V(\hat{\mu}_{MLE})$? Is it consistent?

        \begin{soln}
            \[
                \V(\hat{\mu}_{MLE}) = \V\left(\frac{1}{n} \sum_{i=1}^n x_i\right) = \frac{1}{n^2} \sum_{i=1}^n \V(x_i) = \frac{1}{n^2} \sum_{i=1}^n \sigma^2 = \frac{\sigma^2}{n}
            \]
            Since $\lim_{n \to \infty} \V(\hat{\mu}_{MLE}) = 0$, $\hat{\mu}_{MLE}$ is consistent.
        \end{soln}
    \end{subprob}

    \begin{subprob}
        What is the maximum likelihood estimator $\hat{\sigma}_{MLE}$?

        \begin{soln}
            \begin{align*}
                \hat{\sigma}_{MLE} &= \arg\max_{\sigma} L(\mu, \sigma | x_1, \dots, x_n) \\
                &= \arg\max_{\sigma} \left[ \left(\frac{1}{\sqrt{2\pi\sigma^2}}\right)^n \exp\left(-\frac{1}{2\sigma^2} \sum_{i=1}^n (x_i - \mu)^2\right) \right] \\
                &= \arg\min_{\sigma} \left[ \frac{n}{2}\ln(2\pi) + n\ln(\sigma) + \frac{1}{2\sigma^2} \sum_{i=1}^n (x_i - \mu)^2 \right] \\
                0 &= \frac{\partial}{\partial \sigma} \left[ \frac{n}{2}\ln(2\pi) + n\ln(\sigma) + \frac{1}{2\sigma^2} \sum_{i=1}^n (x_i - \mu)^2 \right] \\
                0 &= \frac{n}{\hat{\sigma}_{MLE}} - \frac{1}{\hat{\sigma}_{MLE}^3} \sum_{i=1}^n (x_i - \mu)^2 \\
                n\hat{\sigma}_{MLE}^2 &= \sum_{i=1}^n (x_i - \mu)^2 \\
                \implies \hat{\sigma}_{MLE} &= \sqrt{\frac{1}{n} \sum_{i=1}^n (x_i - \mu)^2}
            \end{align*}
        \end{soln}
    \end{subprob}
\end{subprobset}
\end{prob}

\newpage

\begin{prob}[\textbf{MLE Basics - Matrix Version}]
Suppose we are given $n$ data points drawn independently from a $d$-dimensional Normal distribution with mean $\mu \in \R^d$ and covariance matrix $\Sigma \in \R^{d \times d}$:

\[ x_1, \dots, x_n \iid \text{Normal}(\mu, \Sigma) \]

\begin{subprobset}
    \begin{subprob}
        What is the likelihood function $L(\mu, \Sigma | x_1, \dots, x_n)$?

        \begin{soln}
            \begin{align*}
                L(\mu, \Sigma | x_1, \dots, x_n) &= \prod_{i=1}^n f(x_i | \mu, \Sigma) \\
                &= \prod_{i=1}^n \frac{1}{\sqrt{(2\pi)^d |\Sigma|}} \exp\left(-\frac{1}{2} (x_i - \mu)^T \Sigma^{-1} (x_i - \mu)\right) \\
                &= \left(\frac{1}{\sqrt{(2\pi)^d |\Sigma|}}\right)^n \exp\left(-\frac{1}{2} \sum_{i=1}^n (x_i - \mu)^T \Sigma^{-1} (x_i - \mu)\right)
            \end{align*}
        \end{soln}
    \end{subprob}

    \begin{subprob}
        What is the maximum likelihood estimate $\hat{\mu}_{MLE}$?

        \begin{soln}
            \begin{align*}
                \hat{\mu}_{MLE} &= \arg\max_{\mu} L(\mu, \Sigma | x_1, \dots, x_n) \\
                &= \arg\max_{\mu} \left[ \left(\frac{1}{\sqrt{(2\pi)^d |\Sigma|}}\right)^n \exp\left(-\frac{1}{2} \sum_{i=1}^n (x_i - \mu)^T \Sigma^{-1} (x_i - \mu)\right) \right] \\
                &= \arg\max_{\mu} \left[ -\frac{1}{2} \sum_{i=1}^n (x_i - \mu)^T \Sigma^{-1} (x_i - \mu) \right]
            \end{align*}
            Since this is a negative sum of quadratic terms, the maximum is achieved when 
            \begin{align*}
                \frac{\partial}{\partial \mu} \left[ -\frac{1}{2} \sum_{i=1}^n (x_i - \mu)^T \Sigma^{-1} (x_i - \mu) \right] &= 0 \\
                \sum_{i=1}^n \Sigma^{-1} (x_i - \mu) &= 0 \\
                \implies \hat{\mu}_{MLE} &= \frac{1}{n} \sum_{i=1}^n x_i \\
                \hat{\mu}_{MLE} &= \bar{x}
            \end{align*}
        \end{soln}
    \end{subprob}

    \begin{subprob}
        What is the expected value of the estimator $\E (\hat{\mu}_{MLE})$? Is it biased?

        \begin{soln}
            \begin{align*}
                \E (\hat{\mu}_{MLE}) &= \E \left(\frac{1}{n} \sum_{i=1}^n x_i\right) = \frac{1}{n} \sum_{i=1}^n \E (x_i) = \frac{1}{n} \sum_{i=1}^n \mu = \mu
            \end{align*}
            Since $\E (\hat{\mu}_{MLE}) = \mu$, $\hat{\mu}_{MLE}$ is unbiased.
        \end{soln}
    \end{subprob}

    \newpage

    \begin{subprob}
        What is the variance of the estimator $\V(\hat{\mu}_{MLE})$? Is it consistent?

        \begin{soln}
            \[
                \V(\hat{\mu}_{MLE}) = \V\left(\frac{1}{n} \sum_{i=1}^n x_i\right) = \frac{1}{n^2} \sum_{i=1}^n \V(x_i) = \frac{1}{n^2} \sum_{i=1}^n \Sigma = \frac{\Sigma}{n}
            \]
            Since $\lim_{n \to \infty} \V(\hat{\mu}_{MLE}) = 0$, $\hat{\mu}_{MLE}$ is consistent.
        \end{soln}
    \end{subprob}

    \begin{subprob}
        \textbf{(Optional, not graded)} What is the maximum likelihood estimator $\hat{\Sigma}_{MLE}$?

        \begin{soln}
            \begin{align*}
                \hat{\Sigma}_{MLE} &= \arg\max_{\Sigma} L(\mu, \Sigma | x_1, \dots, x_n) \\
                &= \arg\max_{\Sigma} \left[ \left(\frac{1}{\sqrt{(2\pi)^d |\Sigma|}}\right)^n \exp\left(-\frac{1}{2} \sum_{i=1}^n (x_i - \mu)^T \Sigma^{-1} (x_i - \mu)\right) \right] \\
                &= \arg\min_{\Sigma} \left[ \frac{n}{2}\ln(2\pi) + n\ln(|\Sigma|) + \frac{1}{2} \sum_{i=1}^n (x_i - \mu)^T \Sigma^{-1} (x_i - \mu) \right] \\
                0 &= \frac{\partial}{\partial \Sigma} \left[ \frac{n}{2}\ln(2\pi) + n\ln(|\Sigma|) + \frac{1}{2} \sum_{i=1}^n (x_i - \mu)^T \Sigma^{-1} (x_i - \mu) \right] \\
                0 &= \frac{n}{2|\hat{\Sigma}|} \frac{\partial}{\partial\Sigma} |\hat{\Sigma}| - \frac{1}{2} \left[ \sum_{i=1}^n (x_i - \mu) (x_i - \mu)^T \right] \hat{\Sigma}^{-1} \frac{\partial}{\partial\Sigma} \hat{\Sigma} \hat{\Sigma}^{-1} \\
                0 &= \frac{n}{2}\hat{\Sigma}^{-1} - \frac{1}{2} \left[ \sum_{i=1}^n (x_i - \mu) (x_i - \mu)^T \right] \hat{\Sigma}^{-1} \hat{\Sigma}^{-1} \\
                \implies \hat\Sigma_{MLE} &= \frac{1}{n} \sum_{i=1}^n (x_i - \mu) (x_i - \mu)^T
            \end{align*}
        \end{soln}
    \end{subprob}
\end{subprobset}
\end{prob}

\newpage

\begin{prob}[\textbf{MLE Basics - Exponential Distribution}]
Suppose we are given $n$ data points drawn independently from an Exponential distribution:

\[ x_1, \dots, x_n \iid \text{Exponential}(\theta) \]

The probability density function for a single point $x_i$ is described as:
\begin{align}
    f (x_i, \theta) =
    \begin{cases}
        \theta \text{exp}(-\theta x_i) & x_i \geq 0 \\
        0 & x_i < 0
    \end{cases}
\end{align}

\begin{subprobset}
    \begin{subprob}
        Describe a real-world outcome that we could model using this distribution.

        \begin{soln}
            This distribution can model things like radioactive decay where $\theta$ is related to the half-life of the particle.
        \end{soln}
    \end{subprob}

    \begin{subprob}
        Derive an expression for $\hat{\theta}_{MLE}$, the maximum likelihood estimator of $\theta$?

        \begin{soln}
            \[
                L(\theta | x_1, \dots, x_n) = \prod_{i=1}^n f(x_i|\theta) = \begin{cases} \theta^n \exp \left(-\theta \sum_{i=1}^n x_i\right) & \forall x_i \geq 0 \\ 0 & \exists x_i < 0 \end{cases}
            \]
            \begin{align*}
                \hat{\theta}_{MLE} &= \arg\max_{\theta} \left[ \theta^n \exp \left(-\theta \sum_{i=1}^n x_i\right) \right] \\
                &= \arg\max_{\theta} \left[ n\ln(\theta) - \theta \sum_{i=1}^n x_i \right] \\
                0 &= \frac{\partial}{\partial \theta} \left[ n\ln(\theta) - \theta \sum_{i=1}^n x_i \right] \\
                0 &= \frac{n}{\theta} - \sum_{i=1}^n x_i \\
                \implies \hat{\theta}_{MLE} &= \frac{n}{\sum_{i=1}^n x_i}
            \end{align*}
        \end{soln}
    \end{subprob}

    \begin{subprob}
        Describe how you could compute $\hat{\theta}_{MLE}$ using only a summary statistic of the data?

        \begin{soln}
            I could use the sample mean of the data points, $\bar{x}$, to estimate $\hat{\theta}_{MLE}$ since $\hat{\theta}_{MLE} = \bar{x}^{-1}$.
        \end{soln}
    \end{subprob}

    \begin{subprob}
        Say you compute the MLE estimator using $n$ data points, and then delete the data. Describe how you can update the MLE estimator with another point $x_{n+1}$? The final equation must be in terms of the summary statistic computed in last part, $x_{n+1}$, and $n$.

        \begin{soln}
            You can use the sample mean of the data points, $\bar{x}$, to estimate the new $\hat{\theta}_{MLE}$ as 
            \[ \hat{\theta}_{MLE} = \frac{n+1}{n\bar{x} + x_{n+1}} \]
        \end{soln}
    \end{subprob}
\end{subprobset}
\end{prob}

\newpage

\begin{prob}[\textbf{MLE without First-Derivatives}]
Suppose we have 5 data points drawn independently from a \textbf{Uniform}(a, b) distribution:

\[ 1.2, 2.1, 1.3, 10.5, 5.0 \]

\begin{subprobset}
    \begin{subprob}
        What are the maximum likelihood estimates for the parameters a and b?

        \begin{soln}
            \[ L(a,b | x_1, \dots, x_n) = \prod_{i=1}^n f(x_i|a,b) = \begin{cases} \left( \frac{1}{b-a} \right)^n & \forall x_i \in [a,b] \\ 0 & \exists x_i \notin [a,b] \end{cases} \]
            To maximize this, you need to find $\arg\min_{a,b} (b - a)$ such that $\forall x_i, x_i \in [a,b]$. This is achieved when $\hat a = \min(x_1, \dots, x_n)$ and $\hat b = \max(x_1, \dots, x_n)$. In this case, the MLE estimates are $\hat a = 1.2$ and $\hat b = 10.5$.
        \end{soln}
    \end{subprob}

    \begin{subprob}
        Suppose that we have $n$ data points $x_1, \dots, x_n$ drawn independently from \textbf{Uniform}(a, b) distribution where $a < b$. What are the maximum likelihood estimators for $a$ and $b$?

        \begin{soln}
            Using the same logic as above, the MLE estimates are $\hat a = \min(x_i)$ and $\hat b = \max(x_i)$.
        \end{soln}
    \end{subprob}
\end{subprobset}
\end{prob}

\begin{prob}[\textbf{Simple Linear Regression via MLE}]
In this problem, we'll use maximum likelihood estimation to fit the "best fit line" between $x$ and $y$ from a dataset of $n$ points $(x_1, y_1), \dots, (x_n, y_n)$.

We assume that each $y_i$ is function of $x_i$ observed with measurement error $\epsilon_i \iid \text{Normal}(0, \sigma^2)$. Formally, we can write:

\begin{align}
    y_i | x_i \sim a x_i + b + \epsilon_i 
\end{align}

\begin{subprobset}
    \begin{subprob}
        For general data $(x_1, y_1), \dots, (x_n, y_n)$, write down the expression for the likelihood and the log-likelihood functions as functions of $a$, $b$, and $\sigma$.

        \begin{soln}
            \begin{align*}
                L(a,b,\sigma | x_1, \dots, x_n, y_1, \dots, y_n) = \prod_{i=1}^n f(y_i|x_i,a,b,\sigma) &= \prod_{i=1}^n \frac{1}{\sqrt{2\pi\sigma^2}} \exp\left(-\frac{(y_i - (a x_i + b))^2}{2\sigma^2}\right) \\
                &= \left( \frac{1}{\sqrt{2\pi\sigma^2}} \right)^n \exp\left(-\frac{1}{2\sigma^2} \sum_{i=1}^n (y_i - (a x_i + b))^2 \right)
            \end{align*}
            \begin{align*}
                \ell(a,b,\sigma | x_1, \dots, x_n, y_1, \dots, y_n) &= \sum_{i=1}^n \ln\left(\frac{1}{\sqrt{2\pi\sigma^2}}\right) - \frac{(y_i - (a x_i + b))^2}{2\sigma^2} \\
                &= -\frac{n}{2} \ln\left(2\pi\right) - n \ln(\sigma) - \frac{1}{2\sigma^2} \sum_{i=1}^n (y_i - (a x_i + b))^2
            \end{align*}
        \end{soln}
    \end{subprob}

    \newpage

    \begin{subprob}
        Derive the maximum likelihood estimates $\hat{a}_{MLE}$ and $\hat{b}_{MLE}$.

        \begin{soln}
            \[ 
                \arg\max_{a, b} \left[ \prod_{i=1}^n \frac{1}{\sqrt{2\pi\sigma^2}} \exp\left(-\frac{(y_i - (a x_i + b))^2}{2\sigma^2}\right) \right] = \arg\max_{a, b} \left[ -\sum_{i=1}^n \frac{(y_i - (a x_i + b))^2}{2\sigma^2} \right] 
            \]
            \begin{align*}
                0 &= \frac{\partial}{\partial b} \left[ -\sum_{i=1}^n \frac{(y_i - (a x_i + b))^2}{2\sigma^2} \right] \\
                0 &= \sum_{i=1}^n (y_i - (\hat{a} x_i + \hat{b})) \\
                0 &= \sum_{i=1}^n y_i - \hat{a} \sum_{i=1}^n x_i - \hat{b}n \\
                \implies \hat{b}_{MLE} &= \frac{1}{n} \sum_{i=1}^n y_i - \hat{a}  \frac{1}{n} \sum_{i=1}^n x_i = \bar{y} - \hat{a} \bar{x} \\
                0 &= \frac{\partial}{\partial a} \left[ -\sum_{i=1}^n \frac{(y_i - (a x_i + b))^2}{2\sigma^2} \right] \\
                0 &= \sum_{i=1}^n (y_i - (\hat{a} x_i + \hat{b})) x_i \\
                0 &= \sum_{i=1}^n y_i x_i - \hat{a} \sum_{i=1}^n x_i^2 - \hat{b} \sum_{i=1}^n x_i \\
                0 &= \sum_{i=1}^n y_i x_i - \hat{a} \sum_{i=1}^n x_i^2 - \left( \frac{1}{n} \sum_{i=1}^n y_i - \hat{a} \frac{1}{n} \sum_{i=1}^n x_i \right) \sum_{i=1}^n x_i \\
                \implies \hat{a}_{MLE} &= \frac{n \displaystyle\sum_{i=1}^n y_i x_i - \displaystyle\sum_{i=1}^n y_i \displaystyle\sum_{i=1}^n x_i}{n \displaystyle\sum_{i=1}^n x_i^2 - \displaystyle\sum_{i=1}^n x_i \displaystyle\sum_{i=1}^n x_i}
            \end{align*}
        \end{soln}
    \end{subprob}

    \newpage

    We will now consider a dataset generated from the model:
    \begin{align}
        y | x \sim \text{Normal}(w_0 + w_1 x, \sigma^2)
    \end{align}

    We assume that the values of $\sigma$ is known and use the dataset to estimate the values of $w_0$ and $w_1$.

    \begin{subprob}
        Show that the given model in (2) and (3) are equivalent.

        \begin{soln}
            The models are equivalent because the likelihood functions are the same. The likelihood function for the model in (2) is:
            \[ L(a,b,\sigma | x_1, \dots, x_n, y_1, \dots, y_n) = \prod_{i=1}^n \frac{1}{\sqrt{2\pi\sigma^2}} \exp\left(-\frac{(y_i - (a x_i + b))^2}{2\sigma^2}\right) \]
            The likelihood function for the model in (3) is:
            \[ L(w_0,w_1,\sigma | x_1, \dots, x_n, y_1, \dots, y_n) = \prod_{i=1}^n \frac{1}{\sqrt{2\pi\sigma^2}} \exp\left(-\frac{(y_i - (w_0 + w_1 x_i))^2}{2\sigma^2}\right) \]
            These are the same if $w_0 = b$ and $w_1 = a$.
        \end{soln}
    \end{subprob}

    We estimate the values of $w_0$ and $w_1$ using the dataset and get $\hat{w_0}$ and $\hat{w_1}$ respectively. We can then predict the outcome using a \textit{fitted model} of the form:

    \[ \hat{f} (x_i) = \hat{y_i} = \hat{w_0} + \hat{w_1} x_i \]

    Given a fitted model, we can measure the \textit{residual sum of squares} as

    \begin{align}
        \text{RSS}( \hat{f} ) = \sum_{i=1}^n (y_i - \hat{f}(x_i))^2 =
        \sum_{i=1}^n (y_i - \hat{w_0} - \hat{w_1} x_i)^2
    \end{align}

    \newpage

    \begin{subprob}
        Show that the RSS is minimized by a fitted model with parameters:
        \[ \hat{w_0} =  \bar{y} - \hat{w_1}  \bar{x}\]
        \[ \hat{w_1} = \frac{\sum_{i=1}^n (x_i - \bar{x})(y_i - \bar{y})}{\sum_{i=1}^n (x_i - \bar{x})^2}\]

        Here, $\bar{x} = \frac{1}{n} \sum_{i=1}^n x_i$ and $\bar{y} = \frac{1}{n} \sum_{i=1}^n y_i$ denote the mean values of $x_i$ and $y_i$ in the dataset.

        \begin{soln}
            Since the RSS is a sum of squares, it is minimized when the gradient is 0.
            \begin{align*}
                0 &= \frac{\partial}{\partial w_0} \left[ \sum_{i=1}^n (y_i - (w_0 + w_1 x_i))^2 \right] \\
                0 &= \frac{\partial}{\partial w_0} \left[ \sum_{i=1}^n (y_i - (\bar{y} - \hat{w_1} \bar{x} + w_1 x_i))^2 \right] \\
                0 &= \sum_{i=1}^n (y_i - \bar{y} + \hat{w_1} \bar{x} - \hat{w_1} x_i) \\
                0 &= \sum_{i=1}^n y_i - \bar{y} n + \hat{w_1} \bar{x} n - \hat{w_1} \sum_{i=1}^n x_i \\
                0 &= \sum_{i=1}^n y_i - \sum_{i=1}^n y_i + \hat{w_1} \sum_{i=1}^n x_i - \hat{w_1} \sum_{i=1}^n x_i = 0 \\
                \implies \hat{w_0} &= \bar{y} - \hat{w_1} \bar{x} \\
                0 &= \frac{\partial}{\partial w_1} \left[ \sum_{i=1}^n (y_i - (\bar{y} - \hat{w_1} \bar{x} + w_1 x_i))^2 \right] \\
                0 &= \sum_{i=1}^n (y_i - \bar{y} + \hat{w_1} (\bar{x} - x_i)) (x_i - \bar{x}) \\
                0 &= \sum_{i=1}^n (x_i - \bar{x})(y_i - \bar{y}) - \hat{w_1} \sum_{i=1}^n (x_i - \bar{x})^2 = 0 \\
                \implies \hat{w_1} &= \frac{\sum_{i=1}^n (x_i - \bar{x})(y_i - \bar{y})}{\sum_{i=1}^n (x_i - \bar{x})^2}
            \end{align*}
        \end{soln}
    \end{subprob}

    \begin{subprob}
        Show that $\V ( \bar{y} ) = \frac{\sigma^2}{n}$. \textit{Note: you can treat the values of $x_1, \dots, x_n$ as fixed in your derivation.}

        \begin{soln}
            Since $y_i \sim \text{Normal}(w_0 + w_1x_i, \sigma^2)$, $\V(y_i) = \sigma^2$.
            \[ \V(\bar{y}) = \V\left(\frac{1}{n} \sum_{i=1}^n y_i\right) = \frac{1}{n^2} \sum_{i=1}^n \V(y_i) = \frac{1}{n^2} \sum_{i=1}^n \sigma^2 = \frac{\sigma^2}{n} \]
        \end{soln}
    \end{subprob}

    \begin{subprob}
        Show that the simple linear regression line will always pass through the point $(\bar{x}, \bar{y})$.

        \begin{soln}
            The equation of the line is $\bar{y} = \hat{w_0} + \hat{w_1} \bar{x}$. Plugging in $\hat{w_0} = \bar{y} - \hat{w_1} \bar{x}$, we get $\bar{y} = \bar{y} - \hat{w_1} \bar{x} + \hat{w_1} \bar{x}$ which becomes $\bar{y} = \bar{y}$ which is true.
        \end{soln}
    \end{subprob}

    \newpage

    We now consider a simpler model where $w_0 = 0$. In this case, we assume that each data point is generated from the model $y | x \sim \text{Normal}(w_1 x, \sigma^2)$, and the fitted model has the form $\hat{f} (x_i) = \hat{w} x_i$.

    \begin{subprob}
        Show that the maximum likelihood estimate can be expressed as 
        \[ \hat{w_1} = \frac{\sum_{i=1}^n x_i y_i}{\sum_{i=1}^n x_i^2} \]

        \begin{soln}
            \[ L(w_1,\sigma | x_1, \dots, x_n, y_1, \dots, y_n) = \prod_{i=1}^n \frac{1}{\sqrt{2\pi\sigma^2}} \exp\left(-\frac{(y_i - w_1 x_i)^2}{2\sigma^2}\right) \]
            \begin{align*}
                \arg\max_{w_1,\sigma} \left[ \prod_{i=1}^n \frac{1}{\sqrt{2\pi\sigma^2}} \exp\left(-\frac{(y_i - w_1 x_i)^2}{2\sigma^2}\right) \right] &= \arg\max_{w_1,\sigma} \left[ -\sum_{i=1}^n \frac{(y_i - w_1 x_i)^2}{2\sigma^2} \right] \\
                0 &= \frac{\partial}{\partial w_1} \left[ -\sum_{i=1}^n \frac{(y_i - w_1 x_i)^2}{2\sigma^2} \right] \\
                0 &= \sum_{i=1}^n (y_i - w_1 x_i) x_i \\
                0 &= \sum_{i=1}^n y_i x_i - w_1 \sum_{i=1}^n x_i^2 \\
                \implies \hat{w_1} &= \frac{\sum_{i=1}^n x_i y_i}{\sum_{i=1}^n x_i^2}
            \end{align*}
        \end{soln}
    \end{subprob}

    \begin{subprob}
        Show that the maximum likelihood estimator is (related to) the correlation between $x$ and $y$.

        \begin{soln}
            The correlation between $x$ and $y$ is given by 
            \[ \rho = \frac{\sum_{i=1}^n (x_i - \bar{x})(y_i - \bar{y})}{\sqrt{\sum_{i=1}^n (x_i - \bar{x})^2 \sum_{i=1}^n (y_i - \bar{y})^2}} \]
            Since the line goes through the origin, we treat it as if the data is centered at the origin, i.e. $\bar{x} = 0$ and $\bar{y} = 0$. Thus,

            \begin{align*}
                \rho &= \frac{\sum_{i=1}^n x_i y_i}{\sqrt{\sum_{i=1}^n x_i^2 \sum_{i=1}^n y_i^2}} \\
                \hat{w_1} &= \rho \sqrt{\frac{\sum_{i=1}^n y_i^2}{\sum_{i=1}^n x_i^2}} \\
                &= \rho \sqrt{\frac{\V(y)}{\V(x)}} \\
                &= \rho \frac{\sigma_y}{\sigma_x}
            \end{align*}
        \end{soln}
    \end{subprob}

    \newpage

    \begin{subprob}
        Show that the best-fit line in this case will pass through the origin.

        \begin{soln}
            \[ \hat{f}(0) = \hat{w_1} \cdot 0 = 0 \]
            Hence, the line passes through the origin.
        \end{soln}
    \end{subprob}

    \begin{subprob}
        Describe a task where we would be better off using the simpler model.

        \begin{soln}
            We would be better off using the simpler model if the relationship between $x$ and $y$ is linear. That is, if you know that when $x = 0$, $y = 0$. One such example is in the case of force and acceleration.
        \end{soln}
    \end{subprob}
\end{subprobset}
\end{prob}

\begin{prob}[\textbf{Complete the Google form at the following \href{https://forms.gle/mHrKRr9banKX15aB6}{link}.}]
    Responding to this form is mandatory as it is worth 1 point. Please note that it's not the content that is graded but whether you filled the form or not.
\end{prob}
\end{probset}
\end{document}